% !TEX program = pdflatex
\documentclass[11pt,a4paper]{article}

% =========================================================================
% FICHA TÉCNICA DE LA PROPUESTA
% Proyecto: CasaScore (AlquilaPuno)
% Gestión de Proyectos TICS (EST315) -- Semestre 2026-II -- Team Sobrevivientes
% =========================================================================

% --- PAQUETES DE CONFIGURACIÓN ---
\usepackage[utf8]{inputenc}
\usepackage[T1]{fontenc}
\usepackage[spanish,es-tabla]{babel}
% Compatibilidad con babel-spanish reciente (flechas de TikZ y signo de porcentaje)
\AtBeginDocument{\shorthandoff{<>}\spanishplainpercent}
\usepackage{geometry}
\geometry{
    top=3.6cm,
    bottom=2.3cm,
    left=2.2cm,
    right=2.2cm,
    % PARTE MODIFICABLE: headheight debe ser mayor que la altura del logo más alto (2.56cm = 73pt).
    headheight=80pt,
    headsep=15pt,
    footskip=20pt
}

\usepackage{graphicx}
\usepackage{fancyhdr}
\usepackage[table]{xcolor}
\usepackage{titlesec}
\usepackage{booktabs}
\usepackage{tabularx}
\usepackage{array}
\usepackage{amsmath}
\usepackage{enumitem}
\usepackage{microtype}
\usepackage{hyperref}

% Logos: funciona a 2 o 3 niveles de profundidad dentro de docs/latex/
\graphicspath{{../../assets/img/}{../../../assets/img/}}

% --- PALETA INSTITUCIONAL FIEI / UNA PUNO ---
\definecolor{AzulUNA}{RGB}{10, 45, 90}
\definecolor{DoradoFIEI}{RGB}{180, 130, 20}
\definecolor{GrisOscuro}{RGB}{50, 50, 50}
\definecolor{AzulClaro}{RGB}{228, 236, 246}

% --- CONFIGURACIÓN DE LISTAS ---
\setlist{topsep=2pt, itemsep=2pt, parsep=1pt}

% --- COLUMNAS DE TABLA ---
\newcolumntype{L}{>{\raggedright\arraybackslash}X}
\newcolumntype{P}[1]{>{\raggedright\arraybackslash}p{#1}}
\newcolumntype{C}[1]{>{\centering\arraybackslash}p{#1}}
\renewcommand{\arraystretch}{1.25}

% --- FORMATO DE ENLACES ---
\hypersetup{
    colorlinks=true,
    linkcolor=AzulUNA,
    urlcolor=AzulUNA,
    citecolor=AzulUNA,
    pdftitle={Ficha Técnica de la Propuesta -- CasaScore (AlquilaPuno)},
    pdfauthor={Team Sobrevivientes}
}

% --- FORMATO DE TÍTULOS ---
\titleformat{\section}{\color{AzulUNA}\normalfont\large\bfseries}{\thesection.}{0.7em}{}[\color{AzulUNA}\titlerule]
\titleformat{\subsection}{\color{AzulUNA}\normalfont\normalsize\bfseries}{\thesubsection.}{0.5em}{}
\titlespacing*{\section}{0pt}{9pt}{4pt}
\titlespacing*{\subsection}{0pt}{6pt}{2pt}

% --- ENCABEZADO Y PIE DE PÁGINA (fancyhdr) ---
\pagestyle{fancy}
\fancyhf{}

% ===== TAMAÑO DE LOS LOGOS (PARTE MODIFICABLE) =====
% logo_fiei.png  : height=1.45cm (lado izquierdo)
% logo_equipo.png: height=2.56cm (lado derecho; 25 % más grande porque la imagen tiene más borde blanco)
% Para cambiar el tamaño, modifique solo el valor height=... de cada \includegraphics.
% Si el logo más alto supera 2.6cm, aumente también headheight y top en la configuración de página.
\fancyhead[L]{%
    \includegraphics[height=1.45cm,keepaspectratio]{logo_fiei.png}%
}
\fancyhead[C]{%
    \scriptsize \textbf{\textcolor{AzulUNA}{UNIVERSIDAD NACIONAL DEL ALTIPLANO DE PUNO}}\\[2pt]
    \tiny \textcolor{GrisOscuro}{FACULTAD DE INGENIERÍA ESTADÍSTICA E INFORMÁTICA}\\[2pt]
    \scriptsize \textbf{\textcolor{AzulUNA}{Gestión de Proyectos TICS (EST315) \textbar\ Team Sobrevivientes}}\\[2pt]
    \footnotesize \textbf{\textcolor{DoradoFIEI}{FICHA TÉCNICA DE LA PROPUESTA}}%
}
\fancyhead[R]{%
    \includegraphics[height=2.56cm,keepaspectratio]{logo_equipo.png}%
}

\renewcommand{\headrulewidth}{0pt}
\renewcommand{\headrule}{}

\fancyfoot[L]{\scriptsize \textcolor{GrisOscuro}{CasaScore (AlquilaPuno)}}
\fancyfoot[C]{\small \thepage}
\fancyfoot[R]{\scriptsize \textcolor{GrisOscuro}{Puno, Perú -- Semestre 2026-II}}
\renewcommand{\footrulewidth}{0.4pt}

\begin{document}

\begin{center}
    {\Large \textbf{\textcolor{AzulUNA}{FICHA TÉCNICA DE LA PROPUESTA}}}\\[3pt]
    {\normalsize \textbf{CasaScore (AlquilaPuno) -- Recomendación inteligente de cuartos y departamentos en alquiler}}\\[4pt]
    \footnotesize \textbf{Fecha:} miércoles 02/09/2026 (semana 2) \quad \textbar\quad \textbf{Versión:} 1.0 \quad \textbar\quad \textbf{Fase:} Inicio
\end{center}

\vspace{0.1cm}

\section{Datos Generales}
\noindent\begin{tabularx}{\textwidth}{@{}P{4.2cm} L@{}}
\toprule
\textbf{Nombre del proyecto:} & CasaScore (nombre de la aplicación: AlquilaPuno) \\
\textbf{Tipo de solución:} & Sistema TIC de recomendación multicriterio (marketplace de alquileres) \\
\textbf{Enfoque principal:} & Motor de recomendación con pesos definidos por el usuario y puntaje explicado de 0 a 100 \\
\textbf{Plataforma:} & Aplicación móvil Android \\
\textbf{Ámbito:} & Ciudad de Puno, entorno de la Universidad Nacional del Altiplano \\
\textbf{Horizonte del proyecto:} & 3 años: MVP en la semana 16 del semestre 2026-II (11/12/2026); año 1 Puno, año 2 Juliaca y año 3 región sur (Cusco y Arequipa) \\
\textbf{Complejidad estimada:} & Media: normalización de variables, modelo multicriterio ponderado y geolocalización \\
\textbf{Enfoque metodológico:} & Cuantitativo y aplicado \\
\textbf{Equipo:} & Team Sobrevivientes \\
\bottomrule
\end{tabularx}

\section{Problema u Oportunidad}
Los estudiantes que llegan a Puno para estudiar en la UNA deben buscar cuarto revisando y comparando a mano anuncios dispersos: precio, distancia a la universidad, seguridad, servicios y reglas de convivencia. No saben si la información es real, están expuestos a estafas y a condiciones que solo descubren después de alquilar. Además, lo que es ``mejor'' cambia según cada persona.

\section{Objetivo General}
Desarrollar un sistema TIC de recomendación multicriterio que permita a los estudiantes buscar, comparar y elegir cuartos y departamentos en alquiler según sus preferencias, presupuesto y ubicación, con un puntaje explicado y confiable.

\section{Objetivos Específicos}
\begin{enumerate}[leftmargin=*]
    \item Registrar las preferencias, el presupuesto y la ubicación de interés del usuario.
    \item Analizar las alternativas disponibles según precio, distancia, seguridad, servicios y condiciones de convivencia.
    \item Generar un ranking personalizado que explique el motivo de cada recomendación.
\end{enumerate}

\section{Beneficiarios}
{\small
\noindent\begin{tabularx}{\textwidth}{@{}P{4.2cm} L@{}}
\toprule
\rowcolor{AzulClaro}
\textbf{Beneficiario} & \textbf{Beneficio} \\
\midrule
Estudiantes (principal) & Encuentran un cuarto adecuado en menos tiempo y con menor riesgo de estafa. \\
Propietarios & Publican su cuarto de forma ordenada y llegan a inquilinos serios. \\
Trabajadores y familias & Usan el mismo sistema para comparar opciones de alquiler. \\
\bottomrule
\end{tabularx}
}

\section{Alcance}
{\small
\noindent\begin{tabularx}{\textwidth}{@{}L L@{}}
\toprule
\rowcolor{AzulClaro}
\textbf{Incluye} & \textbf{No incluye} \\
\midrule
Publicación de cuartos con fotos y datos obligatorios & Cobro de comisiones por el alquiler o intermediación de pagos \\
Configuración de pesos y prioridades del usuario & Firma de contratos de alquiler \\
Cálculo del puntaje CASA SCORE (0 a 100) & Garantía de veracidad absoluta de las publicaciones \\
Ranking personalizado y comparación de opciones & Reemplazo de la visita física al inmueble \\
Mapa de inmuebles respecto a la universidad & Garantía de que el usuario encuentre vivienda \\
\bottomrule
\end{tabularx}
}

\section{Justificación}
Permite encontrar de forma rápida y personalizada el cuarto más adecuado, considerando no solo el precio, sino también la distancia, la seguridad, los servicios y las reglas de convivencia. Usa técnicas de estadística aplicada propias de la formación del equipo (normalización y ponderación multicriterio).

\section{Principal Valor}
Convierte publicaciones dispersas de alquiler en un ranking personalizado y explicado para cada estudiante.

\section{Modelo de Sostenibilidad}
{\small
\noindent\begin{tabularx}{\textwidth}{@{}P{4.2cm} L@{}}
\toprule
\rowcolor{AzulClaro}
\textbf{Fuente de ingreso} & \textbf{Descripción} \\
\midrule
Anuncios recompensados & El usuario elige ver un video corto para desbloquear funciones adicionales. \\
Anunciantes locales & Negocios cercanos a la UNA pagan un espacio publicitario mensual dentro de la aplicación. \\
Suscripción por cuarto & El propietario paga S/ 5 al mes por cada cuarto publicado. \\
Versión Premium & Pago único del estudiante que desbloquea las funciones adicionales sin ver anuncios. \\
\bottomrule
\end{tabularx}
}

\vspace{0.1cm}
\noindent La descarga y el uso básico de la aplicación son gratuitos para el estudiante. La confianza en cada cuarto se muestra con la calificación promedio en estrellas que dan los estudiantes que lo alquilaron; esta calificación no se vende ni modifica el puntaje.

\section{Firmas}
\noindent En señal de conformidad con la presente propuesta, los integrantes del equipo firman a continuación:

\vspace{1.5cm}
\begin{center}
\begin{tabular}{c c c}
    \parbox{6.5cm}{\centering \rule{6.0cm}{0.6pt}\\[5pt] \textbf{Clyde Neil Paricahua Pari}\\[2pt] {\footnotesize Project Manager}}
    & \hspace{0.8cm} &
    \parbox{6.5cm}{\centering \rule{6.0cm}{0.6pt}\\[5pt] \textbf{Ruth Karina Apaza Solis}\\[2pt] {\footnotesize Analista de Planificación}}
    \\[2.0cm]
    \parbox{6.5cm}{\centering \rule{6.0cm}{0.6pt}\\[5pt] \textbf{Eliseo Tarqui Ajahuana}\\[2pt] {\footnotesize Especialista Técnico / Arquitectura}}
    & \hspace{0.8cm} &
    \parbox{6.5cm}{\centering \rule{6.0cm}{0.6pt}\\[5pt] \textbf{Viviana Mary Quisbert Quispe}\\[2pt] {\footnotesize Gestora de Comunicación, Riesgos y Documentación}}
\end{tabular}
\end{center}
\end{document}
